\begin{table}[htbp]
\centering
\caption{P1 系列结果与解释边界。五个 seed 均为算法重复，物理轨迹数始终为三。}
\label{tab:p1-series-summary}
\small
\begin{tabularx}{\textwidth}{@{}l>{\raggedright\arraybackslash}X>{\raggedright\arraybackslash}X@{}}
\toprule
实验 & 数值结果 & 适用范围 \\
\midrule
P1 &
标准 CHNN 的联合场、rollout、线性化误差分别为 \(0.1417\%\)--\(0.2444\%\)、\(0.2131\%\)--\(1.5237\%\)、\(0.3125\%\)--\(0.3286\%\)。 &
仅适用于三条轨迹的 derivative-only 小样本比较。 \\
P1.5 &
residual 相对 physics-only 的场/rollout 平均降幅为 \(90.37\%\)--\(97.29\%\)/\(53.40\%\)--\(95.14\%\)；CHNN 对 gray Real 的场平均降幅为 \(-2.195\%\)。 &
收益有条件且计算更贵；不支持稳定一步场优势。 \\
P1.6 &
\(30/30\) 个约束模型精确通过平衡门，但 Real/CHNN 联合场均值分别恶化至 \(2.872\times10^{-2}\)/\(2.623\times10^{-2}\)。 &
约束满足不等于精度改善。 \\
\bottomrule
\end{tabularx}
\end{table}
